\documentclass[11pt]{article}
\usepackage[margin=1in]{geometry}
\usepackage{amsmath,amssymb}

\title{Direct Langevin Sampling with a Scattering Likelihood}
\author{}
\date{}

\begin{document}
\maketitle

\section{Posterior model}

The observed image is decomposed as
\begin{equation}
    d=s+n,
\end{equation}
where $s$ is the unknown signal and $n=d-s$ is the CIB-like residual.  We
sample the image-space posterior
\begin{equation}
    p(s\mid d)\propto \exp[-U(s)],
    \qquad
    U(s)=U_{\mathrm{like}}(d-s)+U_{\mathrm{prior}}(s).
\end{equation}

For a residual $n$, the STL code computes a scattering-statistic vector
$\phi(n)$.  Each residual is standardized by its own spatial mean and
standard deviation before computing $\phi$, matching the procedure used to
estimate the stored covariance.  With learned mean $\mu_n$ and covariance
$C_n$, the likelihood energy is
\begin{equation}
    U_{\mathrm{like}}(n)
    =\frac{1}{2}\bigl(\phi(n)-\mu_n\bigr)^{\mathsf T}
    C_n^{-1}\bigl(\phi(n)-\mu_n\bigr).
\end{equation}
Only statistics selected by the stored active-statistic mask are used.  A
small diagonal jitter is added to $C_n$ before its Cholesky factorization for
numerical stability.

\section{Fixed power-law signal prior}

The current full run uses a zero-mean stationary Gaussian prior with discrete
power spectrum
\begin{equation}
    P_s(k)=A k^{\alpha}, \qquad k>0,
\end{equation}
where $\alpha=-3.07751$.  The earlier NIFTy slope near $-1.54$ describes the
Fourier amplitude spectrum; the power-spectrum exponent is twice that value.
The normalization $A$ is chosen so that a prior draw
has expected spatial fluctuation RMS $\sigma_s=2.53272$.  These values are
fixed; they are not sampled during the run.  In Fourier space the prior energy
is
\begin{equation}
    U_{\mathrm{prior}}(s)
    =\frac{1}{2}\sum_{k\ne0}\frac{|\widehat{s}(k)|^2}{P_s(k)}.
\end{equation}
The zero Fourier mode is removed because the self-standardized likelihood
cannot identify a constant signal offset.  This fixed-spectrum prior is a
simplification of the earlier hierarchical NIFTy correlated-field prior.

\section{Langevin algorithm and Metropolis correction}

PyTorch automatic differentiation evaluates $\nabla_s U(s)$.  One
unadjusted Langevin (ULA) step is
\begin{equation}
    s_{j+1}=s_j-\epsilon\,\nabla_s U(s_j)
    +\sqrt{2\epsilon}\,\xi_j,
    \qquad \xi_j\sim\mathcal N(0,I).
\end{equation}
The noise and updated image are projected to zero spatial mean.  In ULA this
proposal is always accepted, so finite $\epsilon$ introduces discretization
bias.  The full run instead uses MALA: if $y$ is the proposal from $s$, it is
accepted with probability
\begin{equation}
  \min\!\left[1,
  \frac{e^{-U(y)}q(s\mid y)}{e^{-U(s)}q(y\mid s)}\right],
\end{equation}
where $q$ is the Gaussian Langevin proposal density.  Rejected proposals
leave the state unchanged.  This correction removes the ULA stationary bias,
although it does not by itself guarantee rapid mixing.

The main sampling parameters are:
\begin{description}
    \item[Step size $\epsilon$ (\texttt{--step-size})]
    Controls both the deterministic displacement and the injected noise,
    whose per-step standard deviation is $\sqrt{2\epsilon}$.  A very small
    value mixes slowly; a large value lowers MALA acceptance (or can make ULA
    unstable).  The supplied run uses $\epsilon=10^{-4}$.

    \item[Number of steps (\texttt{--n-steps})]
    Total number of Langevin updates.  Increasing it helps only if the chain is
    stable and has time to forget its initial state.

    \item[Burn-in (\texttt{--burn-in})]
    Number of initial states discarded before saving.  Energy, gradient, and
    residual diagnostics should be approximately stationary after burn-in.

    \item[Thinning (\texttt{--thin})]
    Saves one state every specified number of post-burn-in updates.  It reduces
    file size and correlations between stored states, but cannot repair an
    inadequately mixed chain.

    \item[Initialization (\texttt{--init})]
    The full run first minimizes an STL likelihood plus regularized prior with
    L-BFGS.  This is an STL-based component-separation or MAP-like phase.  Its recovered
    signal is then used as the initial MALA state.  The initializer may use a
    separate, broader power spectrum via
    \texttt{--map-prior-spectral-index}; this affects initialization only, not
    the posterior sampled by MALA.

    \item[Spectral index and RMS]
    \texttt{--prior-spectral-index} controls the relative power across spatial
    scales; more negative values produce smoother fields.
    \texttt{--prior-rms} fixes the overall prior fluctuation scale.

    \item[Covariance jitter]
    \texttt{--covariance-jitter-rel} controls the small diagonal regularizer
    applied to the learned scattering covariance.  It is a numerical
    stabilization parameter, not an additional physical noise term.

    \item[Sampler (\texttt{--sampler})]
    Selects \texttt{ula} or \texttt{mala}.  For MALA, the acceptance rate and
    log acceptance ratio are saved.  Extremely low acceptance means the step
    is too large; very high acceptance with little movement can mean it is too
    small.
\end{description}

\section{Diagnostics}

The run records total, likelihood, and prior energies, the gradient norm and
RMS, and
\begin{equation}
    \rho^2
    =\bigl(\phi(d-s)-\mu_n\bigr)^{\mathsf T}
      C_n^{-1}\bigl(\phi(d-s)-\mu_n\bigr)
    =2U_{\mathrm{like}}.
\end{equation}
For an ideal Gaussian statistic model with $D$ active statistics, a rough
reference is $\rho\sim\sqrt{D}$.  Stable finite values are necessary but not
sufficient: credible posterior samples also require stationary traces and
adequate mixing.  For MALA, the acceptance trace should also be stable and the
posterior summaries should be compared across independent chains.

\end{document}
